\chapter{知识演进与迁移的评价}
\label{chap:ket-evaluation}

花卉识别系统在新地区上线后，平均准确率提高了，却漏掉了更多稀有病害；加入第三个季节的数据后，旧地区的准确率仍高于最初上线时，却低于上一次更新；两个模型合并后，能分别回答物种和病害问题，却不能完成“找出指定物种中患病的花朵”。这些结果并不矛盾。它们考察的是不同对象、不同参照和不同使用条件，不能由一个“迁移效果更好”的总分概括。

第~\ref{chap:ket-reuse-selection}~章已经规定收益的方向，第~\ref{chap:ket-update-maintenance}~至第~\ref{chap:ket-sharing}~章解释了更新、整合、适应与共享的机制。本章要解决的是：怎样把这些方法的目标变成可复查的证据？指标决定要比较什么，数据决定哪些变化真正接受了检验，访问协议和对照决定差异能归于什么，结果分析则限定结论能够外推到哪里。评价中发现的失败可以推动下一轮改进，但不能反过来改写本轮事先约定的成功标准。

本章沿用式~\eqref{eq:ket-reuse-risk-gain}的目标风险$R_T$及收益$\Delta_T=R_T(f_0)-R_T(f_K)$，损失越小越好，因此正值表示有益。准确率等越大越好的分数则用“方法减对照”计算收益。除明确引用的基准设计外，以下花卉数据、三任务矩阵、成本与机构事实均为人工构造的教学例子，不是论文实验结果；其中的初始学习曲线和来源准备成本沿用第~\ref{chap:ket-reuse-selection}~章的例子。

\section{评价目标与指标的确定}
\label{sec:ket-evaluation-metrics}

选择指标前，应把主张写成能被结果否定的句子。例如，“在80个适应标签下，来源初始化降低新地区错误率”不同于“达到90\%准确率所需的全部开发标签更少”；“保留仍有效的旧能力”也不同于“保持旧模型的一切输出”。每项主张都要指定对象、参照、预算与成功条件。多个目标共同存在时，可以给出预定权重，但仍须保留各项原始结果，便于检查权重是否掩盖了必须满足的要求。

\subsection{目标表现与适应效率的度量}
\label{sec:ket-evaluation-target}

固定预算比较要求每个方案在同一预算点留下记录，而不是各取最有利的检查点。记第$r$次运行、预算$b$下的目标分数为$s_r(b)$；$b$可以是标签数、训练步数或设备时间，但一次比较只能使用已经说明的横轴口径。一次记录至少包含数据划分与种子、来源版本、可用目标信息、预算点、正确数或损失和累计资源。使用同一目标样本与随机条件配对比较时，先计算每对运行的收益$d_r(b)$，再汇总
\begin{equation}
  \bar d(b)=\frac1R\sum_{r=1}^{R}d_r(b),\qquad
  s_d(b)=\sqrt{\frac{1}{R-1}
    \sum_{r=1}^{R}\bigl(d_r(b)-\bar d(b)\bigr)^2}.
  \label{eq:ket-evaluation-run-summary}
\end{equation}
这里$R>1$是运行数，$s_d$是运行间样本标准差，不是置信区间。配对差保留了同一批困难样本同时影响两个方案的关系；把两个方案的最优运行相减，会破坏这种对应。独立最终测试还应与调参所用的验证结果分开保存。

回到图~\ref{fig:ket-reuse-learning-curves}。预定适应标签预算为$N=\{20,40,80,160\}$，一次教学运行的迁移准确率依次为$(82,88,91,92)\%$，从头训练为$(65,76,85,91)\%$，每个点都在同一组100张验证照片上计数。这条曲线给出的固定预算收益依次是$(17,12,6,1)$个百分点。横轴只含适应标签，不含来源筛选和验证标签；把全部开发标签作为横轴时，须相应平移或重建每个方案的账目，不能只改轴名。

为说明重复运行的作用，再独立构造两次运行，预算和验证集大小不变。第二次迁移正确数为$(80,87,90,91)$，从头为$(64,75,84,90)$；第三次迁移为$(81,86,89,89)$，从头为$(66,77,85,89)$。在80个标签处，三次配对收益是$(6,6,4)$个百分点，均值为$16/3\approx5.33$个百分点，样本标准差为$\sqrt{4/3}\approx1.15$个百分点。不能只保留第一条迁移曲线的91\%，也不能把三次运行的300次预测当成300张相互独立的新照片：重复测试同一照片所带来的依赖仍然存在。

达到指定效果的代价需要另一份记录。对越大越好的分数、预定阈值$s_\star$和已测预算集合$B$，定义每次运行的经验到达量
\begin{equation}
  b_{\star,r}=\min\{b\in B:s_r(b)\geq s_\star\}.
  \label{eq:ket-evaluation-threshold}
\end{equation}
若集合为空，就记为“截至$b_{\max}$未达到”，而不是令$b_{\star,r}=b_{\max}$。这种记录称为\term{右删失}（right censoring）：只知道到达量超过观察上限，不知道它究竟是多少。阈值采用单个点、连续多个点还是置信下界，必须预先约定；使用验证阈值停止训练以后，还要在独立测试上报告该模型的效果，不能把经验过线解释成总体准确率保证。

在上述三次教学运行中，90\%阈值的迁移到达量为$(80,80,>160)$，从头训练为$(160,160,>160)$。两者截至160的到达比例都是$2/3$，成功运行的条件均值分别为80和160；这两个条件均值不包含第三次运行，不能简称为“全部运行的平均成本”。可以同时报告截至80、160的到达比例，分别为迁移的$2/3,2/3$和从头训练的$0,2/3$，使失败运行仍出现在比较中。若按预算上限截断后平均，迁移得到$(80+80+160)/3\approx106.67$，从头得到160，必须把它称为截断预算统计量；第三个160只是观察终点。它不是未删失的达阈值均值，也不能证明第三次迁移迟早会成功。

取值间隔同样影响结论。第一条曲线只支持“在已测点中，迁移于80过线、从头于160过线”，不支持对40与80之间插值后给出一个精确标签数。预算等比增加时，简单平均四个准确率也不等于按标签投入积分；若计算曲线面积，应写明横轴采用$n$还是$\log n$、积分区间和插值规则。多个任务的准确率、像素误差和推理成功率没有共同量纲，先在任务内计算有方向的收益，再按预定效用尺度和任务权重汇总。

零样本与少样本成绩还必须附带信息条件。零样本可以使用类别名称、属性描述或预训练语料中的知识，但不能一面用目标标签挑提示，一面仍宣称完全没有目标监督；少样本应报告类别数、每类支持样本数、查询集和额外无标签输入。开放世界中，未知检测、建立新类判别规则、旧新类别间校准是三个目标。未知检测需同时检查已知类被拒绝的比例与未知类被接纳的比例；新类学习要在新类留出样本上计分；校准则在旧新类别共同竞争的输出空间中检查。把所有陌生照片都拒绝，可能提高某个拒识指标，却没有学会新增物种。

\subsection{持续能力变化的度量}
\label{sec:ket-evaluation-continuous}

一次更新后的总分看不到能力怎样变化。设按顺序学习$K$个任务，第$t$个阶段结束的模型为$f^{(t)}$，在任务$k$固定评测集$Q_k$上的准确率记为$a_k^{(t)}$。把它们排成矩阵$\bm A$，其中$\bm A_{t,k}=a_k^{(t)}$，行是完成的学习阶段，列是接受评测的任务。$t=0$是尚未经历该序列的共同起点。对角线$a_k^{(k)}$表示刚学完，$k<t$是历史任务区域，$k>t$是尚未学习区域；尚未学习不等于不能评测，缺少输出接口或没有测量才应记缺测。

GEM使用同类矩阵定义最终平均准确率、相对对角线的后向差值，以及相对随机初始化基线的学前零样本差值。本章保留这些比较对象，但把表现记为$a$，避免与风险$R$混淆；若起点已有共同预训练，零样本基线改为这个未经历历史序列的起点，并明确它不再是原论文中的随机初始化\scite{ket-lopezpaz2017-gem}

\subsubsection{历史能力保持与后向迁移的度量}
\label{sec:ket-evaluation-backward}

在阶段$t$，任务等权的\term{阶段平均表现}（stage-wise average performance）只平均已经学习的$t$个任务，即询问当前模型还能完成多少已经承担的工作：
\begin{equation}
  \bar a^{(t)}=\frac1t\sum_{k=1}^{t}a_k^{(t)}.
  \label{eq:ket-evaluation-stage-average}
\end{equation}
这一分母既不是全部计划任务数$K$，也不是训练样本数。任务重要性不同可改用固定权重，并说明在已到达任务上是否重新归一化；若评价所有潜在任务，则另报相应全任务平均。任务集合随$t$扩张，$\bar a^{(t)}$下降也可能是加入了更难的新任务，不能单凭这一条平均曲线判断旧任务遗忘。

保持能力需要逐列比较。对$k<t$，本章定义相对过去已测最好表现的\term{遗忘量}（forgetting）$F_k^{(t)}$，以及相对刚学完时的\term{后向迁移}（backward transfer，BWT）$B_k^{(t)}$：
\begin{equation}
  \begin{aligned}
    F_k^{(t)}&=\max_{k\leq u<t}a_k^{(u)}-a_k^{(t)},\\
    \bar F^{(t)}&=\frac1{t-1}\sum_{k=1}^{t-1}F_k^{(t)},\\
    B_k^{(t)}&=a_k^{(t)}-a_k^{(k)},\\
    \operatorname{BWT}^{(t)}&=\frac1{t-1}\sum_{k=1}^{t-1}B_k^{(t)}.
  \end{aligned}
  \label{eq:ket-evaluation-forgetting-bwt}
\end{equation}
这里$u$枚举已经保存的历史检查点，最大值不含当前阶段。$F>0$表示低于过去峰值，$B>0$表示高于刚学完的表现。若当前创造新高，$F$可以为负；只想统计损失时可另报$\max(0,F)$，但须注明截断，不能无声更换定义。$t=1$时没有历史任务，这两个平均不定义，记为缺测符号而非0。增加历史测量频率还可能捕捉更高峰值，因此比较遗忘量时要匹配检查点。

为说明时间序列，另构造三个两物种识别任务：$\mathcal T_1$是地区甲春季，$\mathcal T_2$是地区甲夏季，$\mathcal T_3$是地区乙秋季。这里的编号仅用于本段序列，不指第~\ref{chap:ket-sharing}~章的物种、病害和定位任务。三个任务的标签含义与输出空间始终相同，各有100张固定、互不混入训练的评测照片，不向模型提供任务身份。图~\ref{fig:ket-evaluation-matrix}~列出人工正确数；因为每列分母都是100，表内数值也就是准确率的百分数。阶段0的三项基线均为50\%；阶段1对任务3没有测量，所以保留缺测，绝不补为零。

\begin{figure*}[htbp]
  \centering
  % Teaching data, not measurements. Every task has 100 fixed evaluation images.
% Rows: (50,50,50), (80,55,missing), (86,82,60), (83,78,88).
\begingroup
\setlength{\tabcolsep}{2.5pt}
\setlength{\fboxrule}{0.6pt}
\renewcommand{\arraystretch}{1.65}
\begin{tikzpicture}[
  x=1mm,y=1mm,
  font=\figurefont\small,
  text=BookInk,
  cell/.style={minimum width=12mm,minimum height=7mm,inner sep=1mm,
    draw=BookMuted,line width=0.6pt},
  flow/.style={-{Latex[length=2mm]},draw=BookMuted,line width=0.8pt}
]
  \node[anchor=north west,inner sep=0pt,font=\figurefont\small\bfseries]
    at (0,4) {(a) 阶段 $\times$ 任务};
  \node[anchor=north west,inner sep=0pt] at (0,-3) {
    \begin{tabular}{@{}p{26mm}*{3}{>{\centering\arraybackslash}p{11mm}}@{}}
      \toprule
      阶段 & $\mathcal T_1$ & $\mathcal T_2$ & $\mathcal T_3$ \\
      \midrule
      0：共同起点 & 50 & 50 & 50 \\
      1：学完$\mathcal T_1$ &
        \fcolorbox{BookTeal}{BookTeal!10}{80} &
        \colorbox{BookAmber!12}{55} &
        \colorbox{BookAmber!12}{\textemdash} \\
      2：学完$\mathcal T_2$ &
        \colorbox{BookBlue!12}{86} &
        \fcolorbox{BookTeal}{BookTeal!10}{82} &
        \colorbox{BookAmber!12}{60} \\
      3：学完$\mathcal T_3$ &
        \colorbox{BookBlue!12}{83} &
        \colorbox{BookBlue!12}{78} &
        \fcolorbox{BookTeal}{BookTeal!10}{88} \\
      \bottomrule
    \end{tabular}
  };
  \node[anchor=west,inner sep=0pt,font=\figurefont\footnotesize]
    at (0,-49) {$k<t$：历史\quad $k=t$：当前\quad $k>t$：未学};
  \node[anchor=west,inner sep=0pt,font=\figurefont\footnotesize]
    at (0,-56) {单位：准确率（\%）；\textemdash：未测};

  \node[anchor=north west,inner sep=0pt,font=\figurefont\small\bfseries]
    at (81,4) {(b) 同列、不同时刻的比较};
  \node[anchor=west,inner sep=0pt] at (81,-6)
    {遗忘：任务1，历史最佳减当前};
  \node[cell,fill=BookBlue!12] (peak) at (88,-16) {86};
  \node[cell,fill=BookBlue!12] (nowf) at (120,-16) {83};
  \draw[flow] (peak.east) -- (nowf.west);
  \node[anchor=west,inner sep=0pt] at (131,-16) {$F_1^{(3)}=3$};

  \node[anchor=west,inner sep=0pt] at (81,-28)
    {BWT：任务1，当前减刚学完};
  \node[cell,draw=BookTeal,fill=BookTeal!10] (learned) at (88,-38) {80};
  \node[cell,fill=BookBlue!12] (nowb) at (120,-38) {83};
  \draw[flow] (learned.east) -- (nowb.west);
  \node[anchor=west,inner sep=0pt] at (131,-38) {$B_1^{(3)}=3$};

  \node[anchor=west,inner sep=0pt] at (81,-50)
    {零样本FWT：任务2，学前减起点};
  \node[cell,fill=BookPaper] (base) at (88,-60) {50};
  \node[cell,fill=BookAmber!12] (before) at (120,-60) {55};
  \draw[flow] (base.east) -- (before.west);
  \node[anchor=west,inner sep=0pt] at (131,-60) {$Z_2=5$};
\end{tikzpicture}
\endgroup

  \caption{持续学习表现矩阵与比较区域，全部数值为人工教学构造。每格为该任务100张固定评测照片上的准确率（\%）；实线框为刚学完的对角线，浅色区域按任务是否已经学习区分，并非数值热度。符号\textemdash 表示未测，不表示准确率为零。右侧按同一列连接参照时刻：遗忘比较过去峰值，BWT比较刚学完，零样本FWT比较学前值与阶段0基线。}
  \label{fig:ket-evaluation-matrix}
\end{figure*}

三个阶段的平均表现分别为$\bar a^{(1)}=80\%$、$\bar a^{(2)}=(86+82)/2=84\%$、$\bar a^{(3)}=(83+78+88)/3=83\%$。在阶段2，唯一历史任务1从80\%升至86\%，所以$F_1^{(2)}=80-86=-6$个百分点，$B_1^{(2)}=86-80=6$个百分点；负遗忘在此只是表示新的最好成绩，不需要把它解释为某种独立机制。

\begin{samepage}
在阶段3，任务1曾达到86\%，现在是83\%，故$F_1^{(3)}=\max(80,86)-83=3$个百分点；任务2刚学完为82\%，现在78\%，故$F_2^{(3)}=82-78=4$个百分点。平均遗忘是$(3+4)/2=3.5$个百分点，分母2只计两个历史任务。相对刚学完的BWT却是
\[
  \operatorname{BWT}^{(3)}
  =\frac{(83-80)+(78-82)}2
  =-0.5\text{个百分点}.
\]

\end{samepage}
任务1相对过去峰值退步了3个百分点，但相对最初学完仍进步3个百分点。两个指标给出不同答案，是因为参照不同，并非计算冲突。一般有
\[
  F_k^{(t)}+B_k^{(t)}
  =\max_{k\leq u<t}a_k^{(u)}-a_k^{(k)}\geq0.
\]
只有刚学完的表现就是此前峰值时，才有$F_k^{(t)}=-B_k^{(t)}$。

这些比较的前提是评测目标仍有效。沿用例~\ref{ex:ket-overview-dated-fact}及第~\ref{sec:ket-update-edit-validation}~节的甲所：2026年7月1日起负责人由甲变为乙，评价应分别保留此前由甲任职的历史问题与此后由乙任职的当前问题，不把纠正过时答案记为有害遗忘。时间条件须写入问题和标准答案，当前事实、历史事实与未变任务分别计分。若物种分类标准改变，也应版本化标准或给出明确映射，再在同一语义下逐列求差。

测试接口也影响“保持”的含义。依靠已知任务身份选取旧任务头，测的是可路由的多头系统；把全部类别放入一个统一输出空间，则还要解决旧新类别竞争。允许调用历史检查点回答旧问题，测的是模型与版本库组成的系统，而不是单个当前模型的保持能力。三者都可以评价，但必须同时报告任务身份、路由正确率、检查点访问与存储条件。

\subsubsection{前向迁移与持续可学性的度量}
\label{sec:ket-evaluation-forward}

历史任务保持得好，并不说明后来的任务学得更快。应分开问：还没有学习新任务时能否直接做得更好，以及开始学习后是否能用更少代价达到要求。\term{零样本前向迁移}（zero-shot forward transfer）衡量前一个问题。令$b_k=a_k^{(0)}$为同一无历史学习起点在任务$k$上的成绩，则
\begin{equation}
  Z_k=a_k^{(k-1)}-b_k,\qquad
  \operatorname{FWT}_0
  =\frac1{K-1}\sum_{k=2}^{K}Z_k.
  \label{eq:ket-evaluation-zero-forward}
\end{equation}
分母为$K-1$，因为任务1之前没有本序列的经验。$Z_k>0$表示学前已经获益。基线和顺序模型必须有相同标签解释、类别描述和预测接口；若新头只能用目标标签拟合，原本的零样本分数就不可得，应记为不适用，不得用拟合后的成绩代替。

图~\ref{fig:ket-evaluation-matrix}~中，任务2的$Z_2=55-50=5$个百分点，任务3的$Z_3=60-50=10$个百分点，所以$\operatorname{FWT}_0=(5+10)/2=7.5$个百分点。阶段1对任务3的缺测不影响这一公式，因为这里只需要每个任务紧邻学习前的单元；若要研究更早阶段对任务3的作用，就必须补测，不能从阶段2的60\%倒推出阶段1的成绩。

CLEAR的原论文补充材料采用另一套汇总：其Backward Transfer是严格下三角的绝对准确率平均，Forward Transfer是严格上三角的绝对准确率平均，都没有减去参照。本例对应的下三角平均为$(86+83+78)/3\approx82.33\%$，不是本章的$-0.5$个百分点BWT；其完整上三角平均因缺少$a_3^{(1)}$而不能计算。CLEAR的Next-Domain Accuracy则只平均紧邻上对角线，本例为$(55+60)/2=57.5\%$，也不是减过50\%基线的7.5个百分点。引用这些名称时，应连同公式、矩阵区域和单位一起给出\scite{ket-lin2021-clear}

学习成本改善需要记录完整适应曲线。令$c_{\star,k}^{\mathrm{hist}}$为带历史经验的模型学习任务$k$、达到预定阈值所需的适应成本，$c_{\star,k}^{0}$为匹配无历史对照的成本。在两者都实际达到阈值、且$c_{\star,k}^{0}>0$时，可定义
\begin{equation}
  E_k=\frac{c_{\star,k}^{0}-c_{\star,k}^{\mathrm{hist}}}
             {c_{\star,k}^{0}}.
  \label{eq:ket-evaluation-learning-cost-gain}
\end{equation}
正值表示节省成本，负值表示更费成本；分母是对照成本，不是两者成本之和。与矩阵配套，另给出一组完整教学预算点$n=(20,40,80,160)$：任务2的历史模型准确率为$(82,90,92,93)\%$，无历史对照为$(70,82,90,92)\%$；任务3的历史模型为$(76,88,90,92)\%$，无历史对照为$(65,76,85,90)\%$。每条曲线使用嵌套支持集、相同每类采样规则和100张固定评测照片，两方案的模型容量与每个预算点的优化预算匹配。90\%阈值下，任务2的到达量为40与80，故$E_2=(80-40)/80=50\%$；任务3为80与160，故$E_3=(160-80)/160=50\%$。

原序列在任务2用了20个标签，在任务3用了40个标签，所以图中相应阶段终点是82\%和88\%。上述扩展曲线从阶段1、2模型的副本出发，只在隔离探测中增加预算，不把探测模型或其评价反馈写回原序列。这样既能复算成本改善，也不会把用于测量的额外训练误算成原系统已经拥有的经验；这些探测开销仍应另行进入评价资源账本。

这组50\%与前面的7.5个百分点没有同一量纲，不能合成一个“前向迁移分数”。如果任务3在160个标签仍未达到阈值，就只能报告$c_{\star,3}^{\mathrm{hist}}>160$及到达比例，不能套入公式制造一个有限节省率。若模型在零成本时已经过线，应直接报告其零样本成功；对照成本也为零时，上式没有定义。

要检验\term{持续可学性}（continual learnability），即经历许多更新后还能否有效吸收新知识，应从多个阶段保存的模型分别出发，学习同一个隔离的后续探测任务。固定新样本、输出头、容量、调参预算和停止条件，比较这些副本的学习曲线，不把它们产生的反馈交还原序列。一个反例是：三个历史阶段对旧任务均保持90\%，但同一探测任务达阈值所需更新步数从40增至80、160，匹配的未经历历史对照始终为80。相应成本改善为$50\%,0,-100\%$，说明“没有忘记”可以与“越来越难学”同时发生。它支持可塑性下降的行为判断，但其原因仍需进一步对照。

若后期模型不断扩容，应加入同等容量的无历史对照，报告增长的存储和计算；若任务顺序不同，应重复顺序而不只重复初始化。长序列还应含环境回归和不同难度任务，并区分历史训练成本与当前探测成本。一个目标适应成功只证明那一次学习可行，不能代替上述跨阶段证据。

\subsection{知识修正与关联一致性的度量}
\label{sec:ket-evaluation-editing}

知识编辑改变的是指定内容及其合理后果。评价应准备四组互相区分的问题：原修改请求、未用于编辑的等义改写或跨语言表达、依赖目标事实的关联推理、按任务语义不应受影响的内容。分别计算修改成功率、表达泛化率、关联正确率和无关内容保持率，分母写成相应请求数或问题数。一次请求有十种改写时，还应先在请求内汇总，防止改写较多的事实得到十倍权重。

单题判分协议必须在运行模型前冻结。封闭式答案可先执行预定的Unicode规范化、大小写与空白处理，再与一组经人工确认的等价答案作精确匹配；日期、数值和单位则使用任务规定的解析器与容差。开放生成若需模型评分器，应固定评分器版本、提示、可见证据和输出标签，并用不参与方法选择的人工标注子集报告一致性；无法解析、拒答和多答案冲突怎样计分也要预先规定。关键词命中不能替代语义正确性，人工事后接受某个有利表述也不能回写评分规则。

给甲所另设一项迁址请求：2026年9月1日起由甲城迁至乙城，两城分别属于甲省和乙省，省份归属不变。这项请求独立于负责人任命。直接问“9月以后所在地是哪座城市”和改写为“9月应到哪座城市拜访”，都应答乙城；问“9月出差应前往哪个省”，则须沿“机构所在地$\to$城市所属省”推出乙省。历史问题“2026年8月应去哪个省”仍应答甲省，另一家未迁址机构的地址也须保持。

在一组教学构造中，10条编辑请求各配1条直接问题、2条未用于编辑的等义改写和1条关联推理题。若直接问题答对9条，20条改写答对16条，10条关联推理题答对6条，则先在每条请求内求正确率、再对请求等权平均，三项结果分别为90\%、80\%和60\%；因为各请求题数相同，这里也等于各组按问题汇总的比例。另有100条无关问题，其中98条输出保持不变，保持率为98\%。直接成功率不能代表关联推理的60\%，输出保持率也不等于正确率：未编辑模型若答错，原样保留仍是错。应在无关集上同时报告编辑前后准确率与预测变化率；必要时再检查概率分布变化，阈值和距离定义事先固定。

\term{局部性}（locality）保护不依赖目标修改的内容，因此迁址后的关联答案改变不计为局部性失败。MQuAKE用多跳问题检验这类后果，并区分反事实编辑与真实时间更新；其有限事实链依赖已知关系与正确前提，不能证明所有推理都一致\scite{ket-zhong2023-mquake}

任务准确率保持、事实保持和行为约束保持也应分开。例如物种分类分数未降，不代表历史负责人问答正确；问答正确，也不代表系统仍遵守不输出未授权记录的约束。需要保护哪些对象，由修改目标及使用要求决定，不应仅挑选容易保持的通用题库。连续编辑时应在每次修改后复查仍有效请求、时间限定的历史请求和互相依赖的请求，记录错误是否累积。

机器遗忘要求移除指定训练记录的影响，参照应是没有使用这些记录的训练过程。设训练集为$S$、撤回记录为$D\subseteq S$，原训练程序为$\mathcal A$，遗忘程序为$\mathcal U$，则比较对象是
\begin{equation}
  f_{\mathrm{unlearn}}=\mathcal U(\mathcal A(S),D),
  \qquad f_{\mathrm{retrain}}=\mathcal A(S\setminus D).
  \label{eq:ket-evaluation-unlearning-reference}
\end{equation}
随机训练下要比较相应模型或输出的分布，而不要求某次运行参数逐位相等。重训练参照应说明初始化、预处理、超参数规则以及所有受撤回数据影响的上游产物；从已经接触撤回记录的检查点继续训练，不自动等于从未见过它们。SISA研究通过限制记录影响范围、重训受影响部分减少撤回成本，其评价同时考察准确率、重训时间和额外存储\scite{bourtoule2021sisa}

经验评价可对照重训练模型检查保留任务效用、撤回样本的预测或生成行为以及预定攻击下的泄漏证据，但撤回样本准确率不必降为零：剩余数据可能仍足以推出相同答案。有限提示下拒答、某次攻击未恢复内容，都不能证明影响已删除。严格删除或近似删除保证的定义、攻击权限及验证口径见第~\ref{sec:privacy-unlearning}~节和第~\ref{sec:privacy-unlearning-validation}~节；本章只把这些要求接入同一实验账目，不用普通任务分数替代隐私要求。

\subsection{多任务共享收益的度量}
\label{sec:ket-evaluation-sharing}

共享是否有益，应以每个任务的独立学习为参照。沿用第~\ref{chap:ket-sharing}~章的物种、病害与定位任务，令$s_k^{\mathrm{joint}}$与$s_k^{\mathrm{ind}}$分别为共同、独立学习的指标，$d_k=1$表示越大越好，$d_k=-1$表示越小越好。需要汇总不同量纲时，预先给定正尺度$c_k$，定义
\begin{equation}
  g_k=d_k\frac{s_k^{\mathrm{joint}}-s_k^{\mathrm{ind}}}{c_k},
  \qquad \bar g=\sum_{k=1}^{K}w_kg_k,\quad
  w_k\geq0,\quad\sum_kw_k=1.
  \label{eq:ket-evaluation-multitask-gain}
\end{equation}
$c_k$只有在被明确规定为政策或效用尺度时，才把各任务的一单位$g_k$置于共同解释下，例如“达到最低服务要求还允许变化多少”。这种尺度及权重$w_k$必须在看结果前确定。取非零独立基线绝对值也可以计算相对基线变化指数，但它不自动建立跨任务共同效用；基线接近零时还会不稳定，因此必须保留原单位结果，且不能把加权平均解释成业务效用。

设独立训练的物种准确率为80\%、病害召回率为70\%、定位平均误差为10像素；共享方案分别为84\%、68\%、9像素。这组教学数值的原始收益是$+4$个百分点、$-2$个百分点和误差减少1像素。若仅为描述相对变化而取$c_k$为各独立基线的绝对值，三项指数为$5\%,-2.86\%,10\%$，等权算术平均约为4.05\%；这个4.05\%只是所选表示下的汇总指数，改变指标表示或尺度就会改变，不能称为跨任务共同效用。最差指数为$-2.86\%$，受损任务比例为$1/3$。平均数为正，病害任务仍受损；如果病害识别是主任务，这个平均收益不能抵消主任务违反要求。

“最差任务”也要说明比较尺度。80\%与10像素不能直接求最小值，可以报告最差归一化收益，或分别报告各任务距其最低服务要求的余量。受损比例则需规定判定规则：按点估计$g_k<0$统计，与要求置信区间完全小于零才记受损，是不同统计量。没有足够证据确认变化方向的任务应保留“不确定”，不能自动归为获益。

\term{帕累托支配}（Pareto dominance）提供不强行合并任务的比较：在已经统一方向的所有目标上，方案A不差于B，且至少一个严格更好，A才支配B。上述共享方案不支配独立方案，因为病害召回率下降。若另一方案为$(84\%,71\%,9\text{像素})$，按点估计它支配$(84\%,68\%,9\text{像素})$；但两者还可能有不同成本。把计算和存储也列为目标后，应重新逐项检查。

比较规则应先说明采用上述不容许任何目标退步的支配关系，还是允许事先规定的小幅退步的实际取舍；后者不等同于上述帕累托支配。有限样本下，应报告配对差及覆盖所比较目标的同时区间，并依据选定规则与零或预设容忍边界比较，以表达证据强度。点估计落在容忍范围内，不等于已经排除更大退步的可能；区间仍跨越判定边界时，应保留不确定结论。训练时找到局部共同下降方向，也不能代替最终测试目标上的这些比较。

相同每任务预算与相同总预算回答不同问题。假设三个独立模型每个花10 GPU小时，总计30；共享训练同样让每个任务参与规定次数的更新，实际总计24 GPU小时。这个比较可考察相同任务暴露下的效果，却不能把共享总投入写成10。若总预算只有10 GPU小时，独立方案须按预定规则把10分给三个模型，共享方案也只能用10；此时比较的是怎样用同一总资源完成整组任务。标注齐全程度、共享骨干计算、缺失标签掩码与调参费用都应列明，任务暴露次数相同不意味着实际算力相同。

\subsection{来源能力保持与组合泛化的度量}
\label{sec:ket-evaluation-composition}

整合后的模型需要接受三种测试：在各来源原任务上保留了多少能力，在联合工作负载上能完成多少请求，在没有完整见过的组合上能否正确计算。设来源$k$在其固定评测集上的原指标为$s_k^{\mathrm{source}}$，整合系统在同一集上的指标为$s_k^{\mathrm{merge}}$，逐来源保持量为$d_k(s_k^{\mathrm{merge}}-s_k^{\mathrm{source}})$。这里参照是对应的来源模型，不是另一个来源中碰巧最好的成绩。除宏平均外，应列出最受损来源及其退化。

例如物种来源为90\%、病害来源为86\%，融合后分别为88\%、87\%，则逐来源变化为$-2,+1$个百分点，任务等权平均从88\%变为87.5\%。如果联合请求中90\%是病害问题，按该工作负载加权却从$0.1\times90+0.9\times86=86.4\%$升至$0.1\times88+0.9\times87=87.1\%$。这两个教学平均都有含义，但只有逐来源记录能显示物种能力下降。提供任务身份选择任务头，与不给身份、在统一标签空间中预测，也必须分别报告。

区分选择与组合需要专门对照。单个最佳来源回答全部问题，是单来源基线；路由器每次选一个来源，检验能否选对；预测集成允许多个来源共同参与，但不必然具备新组合能力；串接模块还需检验中间表示是否被后续计算正确使用。记录来源调用和路由命中率有助于解释过程，却不能仅凭“确实调用了两个来源”就认定完成了组合。

可以在花卉系统中构造属性组合留出。训练覆盖红、黄两种颜色与斑点、条纹两种纹理，每种基本属性都在多种背景、多个对象上出现，但把“红色且有斑点”的完整组合从来源训练、路由训练、提示示例和组合选择集全部留出。最终要求模型找出同时满足两个属性的对象，评价整个合取条件及各属性子判断；同一场景还放入红色条纹、黄色斑点等干扰对象。这样可以排除只看颜色或只看纹理的捷径。若某个强基础模型预训练时已见过相同完整组合，则结论只能称为相对本次任务训练的组合留出，不能宣称所有既有经验中都未见。

保持基本成分覆盖还不够，组合操作也要有学习依据。例如训练中应包含其他属性的合取，才能合理要求把合取规则迁移到新属性对；只提供孤立颜色标签，却要求无提示执行任意逻辑，不是同一个评价问题。完整模式记忆基线、单属性基线、路由基线和联合模块系统应在同一留出组合上比较，并增加新的合取关系或指令顺序，避免一次成功恰好依赖某个词序。参数能相加、模块能连接都只是实现条件，组合留出上的正确行为才是相应证据。

\subsection{全过程资源消耗的度量}
\label{sec:ket-evaluation-resources}

资源账本要沿模型的实际生命周期展开。来源训练和知识准备可能一次发生，筛选和适应可能随任务发生，更新、推理、存储与恢复则随服务时间反复发生。沿用式~\eqref{eq:ket-reuse-amortization}~的口径，规定$C_{\mathrm{prep}}$包含所采用核算范围内的来源训练及准备，$C_{\mathrm{sel},k}$包含全部候选评分与被淘汰试训，$C_{\mathrm{adapt},k}$只计尚未在选择阶段入账的最终适应。设实际服务$m$个任务，服务区间为$[0,H]$；令$\mathcal E_{\mathrm{infer}}(H)$为这段时间实际发生的推理事件集合，$c_e$为事件$e$的增量计算成本。一次共享骨干前向同时产生物种、病害和定位输出时，它是一个事件，而不是三个任务各一次。任务$k$的更新计算合计为$C_{\mathrm{update},k}(H)$。可把计算账目扩为
\begin{equation}
  \begin{aligned}
    C_{\mathrm{total}}(m,H)
    ={}&C_{\mathrm{prep}}+C_{\mathrm{shared}}(H)\\
       &+\sum_{k=1}^{m}\bigl(
        C_{\mathrm{sel},k}+C_{\mathrm{adapt},k}
        +C_{\mathrm{update},k}(H)\\
       &\hspace{24mm}+C_{\mathrm{recover},k}(H)\bigr)
        +\sum_{e\in\mathcal E_{\mathrm{infer}}(H)}c_e.
  \end{aligned}
  \label{eq:ket-evaluation-ledger}
\end{equation}
$C_{\mathrm{shared}}$是未计入各任务的一次或周期性共享维护计算，$C_{\mathrm{recover},k}$包括失败重试与恢复计算；更新成本不能同时再放入共享项。每个推理事件还应记录它服务的任务集合。若管理分析需要把共享事件归因到任务，可另给非负份额$\rho_{e,k}$，并对每个事件满足$\sum_k\rho_{e,k}=1$；任务归因成本$\sum_e\rho_{e,k}c_e$之和仍等于实际推理总成本。每个账目应附事件或版本标识、时间、设备、用途、实际服务任务和是否可复用，不因任务未获益而删除已发生费用。选择中训练的检查点若继续适应，只把后续新增计算计入适应项；重新训练则确实发生了第二份计算。

源模型已经可得时，应同时说明项目增量口径与含来源训练的生命周期口径，不能在同一表中只替某个方案免除历史成本。共享准备费按实际服务任务数$m$和期限$H$摊销，未来可能服务的任务只能做情景分析。教师模型、回放记忆、生成器、特征缓存、索引与版本副本都要登记；把教师在线调用算入训练后，又在同一账目中重复加一次教师计算，同样会失真。

账本不必强行变成一个数。人工标注与审核记录人时，计算记录设备型号、设备时长和必要的运算量，内存报告峰值，持久存储报告权重、数据和缓存的字节数及保留期限。不同时间的存储占用可另记字节时，不能与峰值字节混淆；恢复等待和服务停顿还应报告墙钟时间。参数更少可能降低权重存储，却仍需运行完整骨干和多次检索，所以不能替代实测延迟、吞吐或总计算。训练依据的制作与校验、服务与维护的具体实现分别属于第~\ref{part:efficient-knowledge-use}、\ref{part:systems}~部分的内容范围；当前评价采用本节已经给出的账目、单位与去重规则。

\section{评价数据与基准的选择}
\label{sec:ket-evaluation-data}

指标有了定义，还需要数据让它对应的变化实际发生。一组随机拆分的花卉照片可以评价样本泛化，却没有检验进入新地区；把同一数据按类别拆成多个阶段，可以控制任务干扰，却没有还原真实季节变化。基准的作用是提供可重复的条件，不是给方法颁发适用于所有环境的证明。

\subsection{评价主张与数据覆盖的匹配}
\label{sec:ket-evaluation-coverage}

自然偏移应由采集条件或群体组成的真实变化界定。DomainBed把多领域分类数据、算法实现和模型选择放在统一实验框架中，领域可以由照片风格、数据来源或采集环境区分。在给定数据集内，任务通常仍是同一组类别的识别，改变的是领域而非预测目标。其原论文明确区分训练领域验证、留一领域交叉验证和测试领域验证；最后一种使用目标领域标签作选择，属于更强的先知式条件，不能与完全没有目标数据的领域泛化混作同一比较。DomainBed因此不仅提醒读者留出领域，也要求把如何选超参数算进学习问题；它并不自动检验新类别学习、持续维护或事实修正\scite{gulrajani2021}

WILDS补充了自然应用中的分布变化，如医院、相机位置、地理位置和时间变化，提供对应任务的训练与分布外评价划分。它既包含训练与测试领域分离的领域泛化，也包含领域有重叠但群体比例变化的设定，后者尤其需要检查弱势群体，而不能一概套用整领域留出。与任意施加像素扰动不同，这些划分保留应用任务，改变实际采集条件或群体组成。使用时应说明选中了哪个子数据集、哪个领域或群体、使用哪种官方指标，而不是只报告一个没有组成的“WILDS分数”。自然数据提高了场景真实性，却没有覆盖所有未来偏移，更不能据一次跨医院成功推出跨任务少样本能力\scite{koh2021wilds}

未见任务评价需要隔离任务结构，而不只是隔离照片。Meta-Dataset从多套图像数据构造少样本分类任务，改变每个任务的类别数、支持样本数和类别不平衡程度，并利用部分数据集的类别层级。每个任务取自单一数据集；原始协议把Traffic Signs与MSCOCO完全保留用于评价，其余数据集还做类别划分。因此“在已见数据来源中识别未见类别”与“进入整体未见数据来源”可以分别分析。输入仍以图像分类为主，不能把这个协议直接外推到病斑分割、语言生成或任意新损失的适应\scite{ket-triantafillou2019metadataset}

跨领域少样本还可能改变成像方式。BSCD-FSL从miniImageNet基础类别准备模型，在CropDiseases、EuroSAT、ISIC2018和ChestX上评价，分别涉及植物病害、卫星、皮肤病变和胸部X射线图像。原论文控制了基础模型，并在相同抽样任务上比较元学习与普通微调；其5分类、每类5、20、50个支持样本的协议说明，“少样本”不是固定的一个数字。这里共同保持的是少样本分类形式，改变的是类别与成像领域，不能把数据集间全部差异归结为一个可单独控制的“领域距离”，也没有因此完成长期遗忘评价\scite{ket-guo2020-cross-domain-few-shot}

真实时间流则需要保留时间戳和反馈顺序。CLEAR依据YFCC100M图像的时间信息构造2004至2014年的视觉变化序列，原论文主要固定类别空间考察领域增量。它同时讨论同一时间桶内训练测试划分，以及用当前模型测试下一时间桶的流式协议。后一协议先预测、后揭示标签，随后允许该桶进入训练，因而检验的是面向近未来的表现；这些已经进入训练的数据不再能充当独立的历史保持测试集。要同时研究遗忘，须另留历史审计数据，并说明这已经增加了评价资源。CLEAR没有仅凭真实时间戳就覆盖类别无限增长、无任务边界和事实撤回等所有条件。

连续预训练还要把语料时间、训练阶段和下游探测任务分开。可以在每个语料阶段结束后，用固定下游任务检验能力保持，再用隔离的新任务检验学习成本，同时记录新增词元、旧语料回放及训练计算。只报告新语料损失下降，不能判断下游能力是否保留；只有一次早期自监督预训练，也不能当作持续预训练的长序列证据。CLEAR提供丰富无标签图像，但原论文重点实验中的早期预训练收益，不替代逐时间段持续更新的评价。

内容修改所需的数据围绕事实关系组织。MQuAKE从事实链构造多跳问题，MQuAKE-cf改变部分事实为反事实内容，MQuAKE-t考察真实时间变化。评价保留关系链其余部分，检查被编辑事实的后果是否传播到最终答案；它不同于把同一编辑提示换几个说法。其构造还包含按特定模型可回忆的事实筛选等条件，应记录实体覆盖、事实链长度、编辑条数和问题生成方式。它提供关联推理证据，但不是隐私删除基准，也不穷尽所有局部性与长序列冲突。机器遗忘则需选定真实训练记录与撤回规则，如SISA在Purchase、SVHN等任务中设置的删除请求，并提供匹配重训练参照。随机删除若改成整类或整用户撤回，任务与数据量都可能改变，须另列条件。

多任务共享要求能构造可信的独立对照。对物种、病害与定位任务，应说明哪些照片同时有三类标签，哪些仅有其中一类，缺失标签是否与病害严重度或地区相关。Taskonomy在同一批室内场景图像上提供多种任务标注，通过固定输入像素比较目标任务的变化。原论文重点测量已学来源到目标的迁移关系，不应把其关系分数直接当成联合训练收益。用于共享评价时，仍要重新固定每任务标签、共同对象、训练划分和独立训练预算。若共享模型多看了一批附带物种标签的病害照片，应把增加数据与共享机制分别对照\scite{ket-zamir2018-taskonomy}

组合泛化需把“成分见过、关系未见”落实到划分。SCAN把有限语法生成的导航指令映射为动作序列。随机指令划分、长度外推和把某个原语仅放在单独指令中的划分，要求不同程度的组合泛化；不能用随机划分成绩代表所有系统性组合。它固定了原语含义与解释规则，便于隔离新组合，却没有真实视觉感知、开放语境或无限语法覆盖\scite{ket-lake2018-scan}

COGS把受控英语句子映射为逻辑形式，检查熟悉词语进入新语法角色、熟悉结构的新组合以及更深嵌套等情况。与SCAN的动作序列相比，它扩大了语言结构与语义角色的范围，但仍是规则生成的英语片段，不能自动支持跨语言或视觉组合结论。使用时应分开同分布测试与专门的组合泛化集，并逐类报告失败；把不同泛化类型混成一个均值，会掩盖模型只擅长词汇替换而不擅长结构重组的情形\scite{ket-kim2020-cogs}

\begin{samepage}
模型融合的数据则必须跟随来源定义。Model soups研究同一预训练起点经不同微调配置得到的模型，提供比较最佳单模型、权重平均和预测集成的实验条件；其中同一目标任务上的融合收益不能直接证明不同任务知识已被合并。多任务融合应另列每个来源的原任务测试、共同工作负载与真正组合留出，核对融合系数或来源筛选是否看过它们。模型平均后仍是单个推理模型，说明部署形态没有增加模型副本数，却没有免除训练多个候选和选择配方的成本\scite{ket-wortsman2022-soups}

\end{samepage}

表~\ref{tab:ket-evaluation-claims}~把上述八项作为可多选的评价对象标签，而不是互斥任务类别：上块核对变化与留出，下块核对基准、信息与证据边界，两块行序相同。同一实验可能同时标记多行，但要为每行配套相应划分和访问条件；表中“允许信息”是需要声明的协议，不表示所有同名基准都默认授予这些信息。

\begin{table*}[tbp]
  \centering
  \footnotesize
  \renewcommand{\arraystretch}{1.2}
  \begin{tabularx}{\linewidth}{@{}>{\raggedright\arraybackslash}p{24mm}
    *{2}{>{\raggedright\arraybackslash}X}@{}}
    \toprule
    \multicolumn{3}{@{}l@{}}{\textbf{(a) 评价主张、变化与留出}} \\
    \midrule
    对象标签（可多选） & 所改变因素 & 必要留出 \\
    \midrule
    \mbox{自然偏移} & 领域、采集环境或时期；任务按子集固定 &
      领域、对象组；必要时整段时间 \\
    \addlinespace[2pt]
    \mbox{未见任务} & 类别集合、支持规模、领域 &
      任务与类别；跨领域时留出数据来源 \\
    \addlinespace[2pt]
    \mbox{真实时间流} & 样本随时间到达；可固定标签空间 &
      未来时间段；历史保持另留审计集 \\
    \addlinespace[2pt]
    \mbox{知识编辑} & 指定事实及其依赖结论 &
      请求外改写、事实链、无关与历史问题 \\
    \addlinespace[2pt]
    \mbox{机器遗忘} & 撤回记录、请求粒度与时机 &
      保留测试；撤回审计与攻击测试 \\
    \addlinespace[2pt]
    \mbox{多任务共享} & 任务搭配、标签完整性与共享结构 &
      对象或场景组；每任务独立测试 \\
    \addlinespace[2pt]
    \mbox{组合泛化} & 已有成分的组合关系、深度或次序 &
      完整组合及等价表达；保留成分覆盖 \\
    \addlinespace[2pt]
    \mbox{模型融合} & 来源参数、配方与联合工作负载 &
      逐来源测试、联合测试及组合留出 \\
    \bottomrule
  \end{tabularx}
  \par\medskip
  \begin{tabularx}{\linewidth}{@{}>{\raggedright\arraybackslash}p{24mm}
    >{\raggedright\arraybackslash}p{36mm}
    *{2}{>{\raggedright\arraybackslash}X}@{}}
    \toprule
    \multicolumn{4}{@{}l@{}}{\textbf{(b) 代表基准、信息与证据边界}} \\
    \midrule
    对象标签（可多选） & 代表基准 & 允许信息 & 不能支持的结论 \\
    \midrule
    \mbox{自然偏移} & \mbox{DomainBed}；\mbox{WILDS} &
      来源标签；目标选择信息须另声明 & 任意未来偏移、新任务都可适应 \\
    \addlinespace[2pt]
    \mbox{未见任务} & \mbox{Meta-Dataset}；\mbox{BSCD-FSL} &
      元训练与验证任务；目标支持集 & 少数分类任务代表任意任务分布 \\
    \addlinespace[2pt]
    \mbox{真实时间流} & \mbox{CLEAR} &
      仅已到达数据；反馈后才可训练 & 未来桶正确率就是独立历史保持 \\
    \addlinespace[2pt]
    \mbox{知识编辑} & \mbox{MQuAKE} &
      修改请求及明确允许的辅助事实 & 单跳正确即全局一致或完成删除 \\
    \addlinespace[2pt]
    \mbox{机器遗忘} & \mbox{SISA}中的\mbox{Purchase}、\mbox{SVHN}删除实验 &
      声明剩余数据、检查点与重训参照 & 拒答或一次攻击失败证明严格删除 \\
    \addlinespace[2pt]
    \mbox{多任务共享} & 花卉任务集；\mbox{Taskonomy}数据 &
      匹配标签与样本；缺失标签可见 & 多见数据的收益就是共享机制收益 \\
    \addlinespace[2pt]
    \mbox{组合泛化} & \mbox{SCAN}；\mbox{COGS} &
      已见成分与规则；不泄漏目标组合 & 一类符号组合成功即任意多模态组合 \\
    \addlinespace[2pt]
    \mbox{模型融合} & \mbox{Model soups}；多来源任务测试 &
      模型和选择集；任务头权限另列 & 权重可相加即来源无损或出现新能力 \\
    \bottomrule
  \end{tabularx}
  \caption{评价对象、基准条件与证据边界的非互斥矩阵。上下两块按同一对象标签逐行对应；同一实验可选多行，但每行须分别满足其留出与信息条件。花卉任务集是本章的教学设计，其余名称指正文核对的原始研究；表格不构成跨基准分数排名，机器遗忘列出的删除实验也不能替代独立的删除保证。}
  \label{tab:ket-evaluation-claims}
\end{table*}

\subsection{基准条件与信息来源的核查}
\label{sec:ket-evaluation-provenance}

基准文件没有进入本次训练，并不意味着模型从未接触相关信息。应分别核查预训练中是否见过相同输入、相同输入及标签、完整任务规则、相同组合与近重复样本。只见过某种花的无标签照片，与见过这张测试照片的物种答案，提供的信息不同；见过颜色和纹理两个词，与见过完整“红色且有斑点”指令，也不是同一覆盖。报告应保留这些层次，不把它们压成一个没有依据的“无污染”勾选。

可追溯记录包括来源语料或模型说明、时间范围、版本标识、可获得的数据清单，以及去重的对象与规则。精确哈希可以排除完全相同文件，图像近重复和文本改写还需要相应检索与人工核查；去重阈值、查出的重叠数、排除规则及漏检风险均应说明。只能检查公开子集时，应写“在可访问范围内未发现所查类型的重叠”，不能宣称未知的预训练语料绝无测试内容。

运行时的信息也要逐项登记：是否能访问旧原始数据、仅保留记忆还是保存全部检查点；是否提供任务身份；标签在预测前还是预测后到达；是否允许目标无标签批量联合处理。一个带任务头的分类分数和一个不知任务身份的分类分数，所测难度不同。来源筛选、提示示例、校准阈值和人工失败分析用到的目标数据也都属于访问记录，不能因其未进入梯度就忽略。

核查产出应能回答一条具体结论的依据。例如，“新地区结果使用了该地区30张带标签选择照片、80张适应照片和未标注校准批次，未用最终测试标签选择来源”，比“采用标准迁移协议”更可复查。哪些对象必须整体留出、哪些数据可共享，由第~\ref{sec:ket-evaluation-splits}~节的划分协议落实；若发现重叠，应重做独立评价或收窄结论，而不是仅在脚注承认后继续使用原来的强主张。

\subsection{补充场景与失败样本的构造}
\label{sec:ket-evaluation-stress}

已有基准没有覆盖的条件，应由研究目标决定是否补测。花卉系统若要跨季节长期使用，补充场景应包括稀有病害、遮挡、进入过的旧地区再次出现，以及比原基准更长的任务序列；编辑系统若要持续维护，则需包含同一事实的再次更新、相互冲突的请求和缺少依赖前提的问题。样本的纳入规则应在看结果前说明；事后挑出的失败可以作为诊断集，但不能再估计预先未定义的总体失败率。

受控构造应尽量只改变待检验因素。比较组合能力时，保持颜色、纹理、对象数量和背景分布相同，只改变属性关系是否在训练中共同出现；比较分布适应时，固定目标标签和样本量，改变拍摄光照；比较关联编辑时，固定目标事实、语言模板及编辑条数，只改变依赖链长度。链更长若同时包含更多冷门实体，就不能把失败全部归给推理深度，应补做实体熟悉度匹配。

构造还必须满足任务语义。若一个机构的迁址请求没有生效日期，历史和当前问题就可能没有唯一正确答案；应补齐日期，或把“信息不足时明确说明无法确定”作为目标。冲突请求则须规定版本优先级或待核验状态，不能期待模型凭空知道哪条更可信。病害合成若违反真实成像和病理条件，只能检验模型对所定义图案的响应，不能伪装成临床或田间证据。

教学构造用于复算定义，受控实验用于区分候选解释，真实数据用于检验部署条件。本章的三任务矩阵属于第一类，不能给它加上随机误差棒就变成实验。正式补充数据应保存生成规则、随机种子、采样范围、质量核查和失败样本的完整纳入清单；所有方法使用同一套条件，并同时呈现成功、失败和无法评价项。

\section{对照实验与评价协议的设计}
\label{sec:ket-evaluation-protocol}

数据确定以后，还要规定谁在何时能看到哪些信息、哪些方案接受相同对照，以及随机重复覆盖什么变化。这些选择应在最终比较前固定。若评价中途改变了选择规则或成功阈值，应把新结果标为探索性分析，并为确认性结论安排新的独立数据。

\subsection{数据划分与访问协议的设计}
\label{sec:ket-evaluation-splits}

同一份花卉数据可以按不同层级留出。样本层留出要求训练和测试照片不重叠，并把同一植株的连拍、裁剪及增强视为同一对象组；类别层留出要求最终物种未进入训练类别；任务层留出要求新的分类问题或预测目标未参与元训练与选择；领域层留出隔离地区、设备或采集来源；时间层留出只允许过去信息预测以后。它们不能互相替代：照片互不相同，仍可能共享同一植株；物种不同，仍可能来自相同地区；任务支持集与查询集分开，也可能全都属于元训练见过的任务族。

例如把每个地区每种花的照片随机拆开，只检验同一组类别与领域中的新样本。把物种丙、丁整体留下，检验未见类别，但若仍从已见地区拍摄，就没有同时检验领域外推。再把地区乙整体留下，才增加领域留出；按2026年7月划分过去与未来，又增加时间条件。联合留出往往更难，也可能使数据不足，应写清哪些层级被同时控制，而不是笼统使用“严格划分”。

图~\ref{fig:ket-evaluation-boundaries}~给出选择、适应和最终评价的职责。来源数据形成候选知识，目标选择集用于决定来源、超参数与停止规则，适应集拟合最终任务模型，最终评测集只用于锁定方案后的只读计分。少样本任务可通过开发任务选择适应规则，在最终任务仅使用其支持集；若规定目标选择集允许并入最终适应集，这两个用途就不再独立，应明确标注共享，按唯一样本计标注成本，最终评测集仍留在外部。

图中底部的三个网格使用完全相同的九张照片。样本留出后，训练仍覆盖A、B、C三个物种和两个地区；按类别留下整行C，才要求识别未见物种；按领域留下整列3，则要求在未见地区乙识别已经见过的物种。样本没有重复，只满足了这些要求中最基础的一项。

元学习沿用第~\ref{sec:ket-transfer-task-splits}~节的$\mathcal T_k\sim\Pi$、支持集$S_k$与查询集$Q_k$。元训练查询标签参与学习规则更新，所以不是独立最终证据；元验证任务用于比较候选规则，选择完成后锁定规则，才将它交给最终测试任务使用。交接的是规则，不是元验证任务或它的标签；最终任务的$S_k$用于允许的适应，$Q_k$的标签仅用于计分。一个任务抽样器在已见类别上产生了新的一组支持照片，并不足以证明类别或领域留出。还应声明是否允许共同查看同一任务的全部查询输入；允许这种联合处理的直推式协议，与每次只接收一个查询的归纳式协议具有不同信息条件。

\begin{figure*}[htbp]
  \centering
  % Solid: fitting/update or labeled rule handoff; dashed: selection evidence.
% Dotted: scoring. Task collections never flow into one another.
\begin{tikzpicture}[
  x=1mm,y=1mm,font=\figurefont\footnotesize,text=BookInk,
  box/.style={draw=BookMuted,line width=0.7pt,minimum height=13mm,
    text width=27mm,align=center,inner sep=1.5mm},
  train/.style={-{Latex[length=2mm]},draw=BookTeal,line width=0.9pt},
  select/.style={-{Latex[length=2mm]},draw=BookAmber,line width=0.8pt,dashed},
  score/.style={-{Latex[length=2mm]},draw=BookBlue,line width=0.8pt,dotted},
  title/.style={anchor=west,inner sep=0pt,font=\figurefont\small\bfseries}
]
  \path[use as bounding box] (0,-178) rectangle (169,17);
  \node[title] at (0,13) {(a) 开发与独立最终评价};
  \node[box,fill=FigureInput!10] (source) at (16,0) {来源数据\\形成候选};
  \node[box,fill=FigureModel!10] (knowledge) at (60,0) {候选知识\\确定复用配置};
  \node[box,fill=FigureModel!10] (adapt) at (105,0) {目标适应\\执行已选规则};
  \node[box,fill=FigureOutput!10] (frozen) at (149,0) {锁定模型\\封存预测};
  \draw[train] (source.east) -- (knowledge.west);
  \draw[train] (knowledge.east) -- (adapt.west);
  \draw[train] (adapt.east) -- (frozen.west);

  \node[box,fill=FigureInput!10] (selection) at (60,-26)
    {目标选择集\\来源、超参数、停止};
  \node[box,fill=FigureInput!10] (support) at (105,-26)
    {目标适应集\\拟合任务模型};
  \node[box,fill=FigureInput!10] (test) at (149,-26)
    {最终评测集\\只读计分};
  \draw[select] (selection.north) -- (knowledge.south);
  \draw[train] (support.north) -- (adapt.south);
  \draw[score] (frozen.south) -- (test.north);
  \draw[draw=BookMuted,line width=0.6pt]
    (selection.south) -- (60,-38) -- (105,-38) -- (support.south);
  \node[anchor=north,inner sep=1mm] at (82.5,-38)
    {仅按预定协议共享};
  \node[align=left,anchor=west,inner sep=0pt,text width=30mm] at (0,-27)
    {最终结果不回流\\选择与适应};

  \node[title] at (0,-54) {(b) 元学习：任务层留出};
  \node[box,text width=31mm,fill=FigureInput!10] (metatrain) at (18,-68)
    {元训练任务\\$S_k$拟合，$Q_k$更新};
  \node[box,text width=31mm,fill=FigureModel!10] (candidates) at (62,-68)
    {候选规则};
  \node[box,text width=31mm,fill=FigureInput!10] (metaval) at (18,-88)
    {元验证任务\\试用候选规则};
  \node[box,text width=31mm,fill=FigureModel!10] (lockedrule) at (62,-88)
    {锁定规则};
  \node[box,text width=31mm,fill=FigureOutput!10] (metatest) at (62,-111)
    {最终测试任务\\$S_k$适应，$Q_k$计分};
  \draw[train] (metatrain.east) --
    node[midway,above,inner sep=1mm] {学习} (candidates.west);
  \draw[select] (candidates.south) --
    node[midway,right,inner sep=1mm] {选择} (lockedrule.north);
  \draw[select] (metaval.east) --
    node[midway,above,inner sep=1mm] {验证} (lockedrule.west);
  \draw[train] (lockedrule.south) --
    node[midway,right,inner sep=1mm] {交接规则} (metatest.north);
  \node[anchor=west,align=left,text width=33mm,inner sep=0pt] at (0,-111)
    {任务集合分开\\未见类别另做\\类别层留出};

  \node[title] at (91,-54) {(c) 连续学习：时间顺序};
  \node[box,text width=56mm,minimum height=9mm,fill=FigureOutput!10]
    (predict) at (128,-68) {状态$\sigma^{(t)}$：先预测并封存};
  \node[box,text width=56mm,minimum height=9mm,fill=FigureInput!10]
    (feedback) at (128,-85) {反馈到达：计分并按协议更新};
  \node[box,text width=56mm,minimum height=9mm,fill=FigureModel!10]
    (next) at (128,-102) {状态$\sigma^{(t+1)}$：仅服务未来};
  \draw[score] (predict.south) -- (feedback.north);
  \draw[train] (feedback.south) -- (next.north);
  \node[anchor=west,inner sep=0pt] at (91,-115)
    {独立运行重置；流内状态按协议保留};

  \node[title] at (0,-132) {(d) 同一记录池，不同留出层级};
  \foreach \shift/\kind/\heading in {0/1/样本留出,55/2/类别留出,110/3/领域留出} {
    \node[anchor=west,inner sep=0pt] at ({\shift+10},-140) {\heading};
    \foreach \r/\letter in {1/A,2/B,3/C} {
      \foreach \c in {1,2,3} {
        \ifnum\kind=1
          \ifnum\r=\c
            \fill[BookAmber!18]
              ({\shift+10+7*(\c-1)},{-146-7*(\r-1)})
              rectangle ({\shift+17+7*(\c-1)},{-153-7*(\r-1)});
          \fi
        \fi
        \ifnum\kind=2
          \ifnum\r=3
            \fill[BookAmber!18]
              ({\shift+10+7*(\c-1)},{-146-7*(\r-1)})
              rectangle ({\shift+17+7*(\c-1)},{-153-7*(\r-1)});
          \fi
        \fi
        \ifnum\kind=3
          \ifnum\c=3
            \fill[BookAmber!18]
              ({\shift+10+7*(\c-1)},{-146-7*(\r-1)})
              rectangle ({\shift+17+7*(\c-1)},{-153-7*(\r-1)});
          \fi
        \fi
        \draw[BookMuted,line width=0.5pt]
          ({\shift+10+7*(\c-1)},{-146-7*(\r-1)})
          rectangle ({\shift+17+7*(\c-1)},{-153-7*(\r-1)});
        \node[inner sep=0pt] at ({\shift+13.5+7*(\c-1)},{-149.5-7*(\r-1)})
          {\letter\c};
      }
    }
  }
  \node[anchor=west,inner sep=0pt] at (0,-174)
    {A、B、C为物种；第3列来自地区乙，其余来自地区甲；阴影为留出。};
\end{tikzpicture}

  \caption{选择、适应与最终评价的信息边界。实线箭头表示拟合、按协议更新或已标注的规则交接，虚线表示选择依据，点线表示只读计分。元验证任务帮助选择候选规则，锁定后交给最终任务的是规则；三个任务集不沿箭头流转，最终查询标签只计分。选择集与适应集仅在预先声明时可共享，最终评测标签与结果不得返回本轮开发。底部用同一组九张教学照片示意样本、类别和领域留出，阴影格为留出记录。在线反馈只服务未来，不能重算过去预测。}
  \label{fig:ket-evaluation-boundaries}
\end{figure*}
\FloatBarrier

测试时适应尤其需要写清状态。设预测器的可变状态为$\sigma^{(t)}$，其中可能包含参数、归一化统计、优化器、记忆、缓存和随机数状态。独立任务或独立运行开始时，应恢复约定起点；同一任务内是否跨样本保留状态、结束后是否丢弃，都要固定。只恢复参数却留下上一任务的归一化统计或缓存，并没有完成协议要求的重置。

在线评价通常采用\term{先预测后更新}（predict-then-update）：用当前状态产生$\hat y_i$并封存，真实标签到达后才计分和更新，更新后的状态服务未来样本。若标签延迟$d_i$步，时刻$i$只能使用已经到达的旧反馈，不能提前借用$y_i$。另一种测试时适应协议允许先用当前无标签输入更新、再预测同一输入，必须说明更新使用了哪些无标签样本、是否读取整个测试批量；使用真实标签更新后再在同一对象上计分则是训练拟合，不能冒充预测成绩。流式标签进入后续训练属于预定在线协议，并不意味着可以事后改变已经记录的预测。

持续学习可分别使用真实时间顺序、受控任务排列和环境回归。真实时间顺序保持因果访问，不应随意打乱以获得更好的均值；受控任务排列用于隔离顺序效应，可重复随机排列；环境回归检验离开后再次遇到旧条件时能否恢复与继续学习。三者应分别报告。若要在在线训练之外保留独立最终审计，审计集不得进入上述反馈流。

知识编辑的划分要以请求及关系族为单位。原请求可用于修改，等义改写、跨语言表达、关联推理和无关任务须分别保留；直接答案的重复模板不能同时算作多个独立事实成功。评价多跳能力时，编辑请求可以提供被修改事实，但不能顺便把待测多跳问答的完整答案作为示例交给模型。比较访问条件时，可以固定模型与数据，另设有无源数据、有无任务身份、是否保留测试状态的实验组；这些组检验信息权限的作用，不应把强权限组的成绩当成同权限算法改进。

\subsection{比较基线与消融实验的设计}
\label{sec:ket-evaluation-baselines}

基线应回答研究主张中最直接的替代解释。检验来源初始化是否有益，比较相同结构、相同目标数据下的从头训练与普通微调，并加入冻结来源表示和冻结随机表示，区分已有表示与目标重建；检验持续保持，比较朴素顺序训练与同样记忆预算的简单回放；检验共享，比较每个任务独立训练及相同总预算下的独立方案。基线也要有合理调参机会，不能以未收敛的随机初始化作唯一对照，然后把所有差异归给新机制。

若主张来自专用梯度协调器，单任务对照还不够，因为它同时取消了共享结构。应固定骨干、任务头、数据、损失单位、训练起点、基础优化器与停止规则，加入等权联合训练，再在预定搜索预算内扫描固定损失权重或任务采样概率。各方案预先报告候选范围、验证选择规则和同口径的搜索预算，最终测试只评价锁定方案；逐任务梯度、协调求解及被淘汰配置的费用均须入账。比较应保留逐任务性能取舍，而非只挑一个有利平均。Xin等在其语言任务与已测方法上发现，专用方法没有超出简单采样扫描得到的性能取舍；这支持把充分调参的简单联合训练列为直接对照，不是否定所有协调器，也不涵盖该论文未比较的Nash-MTL\scite{ket-xin2022-mto}

参数融合应至少比较按选择集选出的最佳单来源、直接预测集成与融合模型。最佳来源不能按最终测试逐任务挑选而不声明先知条件；多头集成得到的访问优势也要保留。若融合比集成差但占用更少推理资源，这可能是有价值的取舍，应报告效用与成本，不能仅挑单模型基线宣称在所有意义上更优。组合模块还需加入只路由单来源、移除某一模块或替换为无信息输入的对照，检查所需的多个成分是否真正贡献。

元学习的对照应固定预训练来源、骨干容量和目标支持集，再比较普通适应、只用初始化不适应、完整元学习适应。若普通适应没有得到同样预训练，就无法把差异归给学习规则。对来源选择，可用固定默认来源、随机选来源、代理评分选择和等预算目标试训比较；同时保留搜索开销。用最终测试选出的最佳来源可作诊断参照，但不能作为可部署选择器的成绩。

\term{消融实验}（ablation study）通过移除或替换一个部件，检查完整方案的变化是否符合机制解释。消融回放时，应说明同时减少了旧数据访问还是仅替换记忆内容；消融共享结构时，应报告参数量和每任务计算是否改变；移除适应步骤时，可另设把节省预算用于普通更新的对照，以区分专门机制与更多更新的作用。没有一种消融天然只改变一个因素，实际改变的条件要逐项登记。

同时访问全部历史数据的联合训练可以作为强参照，帮助判断受限记忆付出了多少代价。但它违反了很多持续学习协议的数据访问限制，不能与受限方法混写排名。有限计算、固定网络和某次优化得到的联合训练成绩，也不是严格数学上界：另一个优化过程可能超过它。表述为“本配置下的全历史联合训练参考”更准确。

\subsection{资源预算与重复实验的设计}
\label{sec:ket-evaluation-repeats}

固定总预算实验要在预算内同时支付来源选择、失败试训和最终适应；达到指定效果实验则记录整个过线过程和未过线运行。来源准备是否计入、共享成本按多少任务摊销、服务请求量和截止时间，都应在比较前确定。预算不允许完成某方案时，应报告不可执行或未完成，而不是只统计成功完成的便宜运行。

资源匹配需说清控制了哪个量。标签数相同，不代表人工成本相同；步数相同，不代表批量大小、网络规模或计算相同；GPU小时相同，不代表硬件与利用率相同；权重文件相同大小，也不代表峰值内存相同。可以分别设计标签预算、实际计算预算和内存约束实验，不必假装所有量都能同时相等。无法匹配的差异应保留在结果表中，并限制相应的机制结论。

要从重复记录得到区间，须先确定估计对象与独立单位。沿用第~\ref{sec:ket-evaluation-target}~节在80个标签处的配对收益$(6,6,4)$个百分点，另设一项只演示区间计算的教学协议：固定任务、训练与支持数据、一组预先封存的测试照片、模型配置及超参数，只重新抽取运行随机性。每对运行让两方案共享能够匹配的随机条件，各对独立抽取种子。因此独立单位是一对完整运行，不是其中的一张照片；这里也不把前文用于选模型的验证成绩改称独立测试证据。

记$d_r$为第$r$对运行的准确率差，估计对象$\mu_d=\mathbb E_{\mathrm{run}}[d_r]$是在上述固定条件下、对运行随机性取平均的收益。为演示小样本区间，假设这些运行差独立同分布且服从正态分布。沿用式~\eqref{eq:ket-evaluation-run-summary}~的$\bar d,s_d$，均值标准误的估计为$\widehat{\operatorname{SE}}(\bar d)=s_d/\sqrt R$，即均值因有限运行数产生的波动尺度。预定置信水平为95\%，以$t_{\nu,p}$表示自由度$\nu$的Student $t$分布的$p$分位数，区间为
\begin{equation}
  \bar d\ \pm\ t_{R-1,0.975}\,
    \frac{s_d}{\sqrt R}.
  \label{eq:ket-evaluation-paired-interval}
\end{equation}
本例$R=3$，自由度为2，代入得到
\[
  \begin{aligned}
    \bar d&=\frac{16}{3},\qquad s_d=\sqrt{\frac43},\\
    \widehat{\operatorname{SE}}(\bar d)&=\frac{\sqrt{4/3}}{\sqrt3}
      =\frac23,\\
    \bar d\pm t_{2,0.975}\frac23
      &\approx5.333\pm4.303\times0.667\\
      &\approx[2.46,\,8.20]\text{个百分点}.
  \end{aligned}
\]
95\%描述区间构造规则的长期覆盖率：若上述假设成立，反复抽取成组运行并计算区间，约95\%的区间会覆盖固定的$\mu_d$；不是说某次运行有95\%概率获益。三次人工数据不能验证独立同分布或正态假设，也不是推荐的实际运行次数。这一区间不覆盖新照片、新任务、重新抽取支持集或重新调参所产生的不确定性。基准比较需记录这些随机来源；均值差和“两方案谁胜出的概率”是不同统计量，不能直接套用后者的区间公式\scite{ket-bouthillier2021-variance}

重复实验至少要分清固定任务内的运行随机性与任务本身的变化。设$g_{k,r}$是任务$k$上第$r$次配对运行的收益，每任务运行$R$次，$K$个评价任务等权。可以分别报告
\begin{equation}
  \begin{aligned}
    \bar g_k&=\frac1R\sum_{r=1}^{R}g_{k,r},&
    \bar g&=\frac1K\sum_{k=1}^{K}\bar g_k,\\
    s_{\mathrm{run}}^2
      &=\frac1K\sum_{k=1}^{K}
        \frac{\sum_{r=1}^{R}(g_{k,r}-\bar g_k)^2}{R-1},\\
    s_{\mathrm{task}}^2
      &=\frac{\sum_{k=1}^{K}(\bar g_k-\bar g)^2}{K-1}.
  \end{aligned}
  \label{eq:ket-evaluation-two-variations}
\end{equation}
这里$R,K>1$。前者概括同一任务内运行的波动，后者概括任务均值之间的差异；后者还含有限重复估计造成的噪声，不能不经建模就称为纯粹的任务总体方差。这两个描述量都不是标准误，不能直接拿来给$\bar g$画置信区间。

考虑两个任务各有三次收益，分别为$(5,6,7)$和$(-3,-2,-1)$个百分点。任务均值为6和$-2$，总均值为2；任务内样本方差均为1，故$s_{\mathrm{run}}=1$个百分点；任务均值之间的样本方差为32，故$s_{\mathrm{task}}=\sqrt{32}\approx5.66$个百分点。多跑几次可以更准确地估计每个任务的均值，却不能把只有两个任务变成对广泛任务分布的充分覆盖。平均正2也没有取消第二个任务的负收益。

不采用正态运行差模型时，可以用\term{配对自助法}（paired bootstrap）近似区间：反复有放回地抽取实际独立单位，两方案始终共用抽样索引，再观察所需统计量怎样变化。以固定任务、固定测试集的运行均值为例，可按下面的步骤执行。
\begin{enumerate}
  \item 保存每对运行中两个方案的分数及共同随机条件。预先固定统计量、任务权重、95\%置信水平、重抽样次数$L$和随机种子；例如取$L=10000$，它是计算次数，不是新增实验次数。
  \item 每次从$R$个配对运行编号中有放回地抽取$R$个索引。两方案使用同一串索引，逐对重算分数差，再求均值。对多个固定任务仍按原权重汇总，不把任务多跑几次变成更高权重。
  \item 重复$L$次，保存每次的均值差，取其经验2.5\%与97.5\%分位数作为百分位自助区间。报告原始运行数、抽样单位、分位数规则和$L$；区间覆盖率是近似的，只有三对记录时重抽样不能补足分布信息。
\end{enumerate}

多任务下，抽样索引须遵守共同来源造成的相关。若同一个元训练种子产生的模型服务全部$K$个任务，就把该种子下两方案的整列任务记录作为一组：抽中运行$r$，须同时保留$(g_{1,r},\ldots,g_{K,r})$，不能逐任务另抽元训练种子。只有各任务确实独立训练时，才可在任务内分别抽运行。若只研究固定模型对新测试对象的平均收益，则改为对独立对象组抽样，两方案共用组索引，组内照片全部保留；这得到的是测试对象抽样区间，不包含重训练的不确定性。

固定任务集的区间不能支持对$\Pi$中新任务的外推。若评价任务确实按$\Pi$独立同分布抽取，且结论只针对一对已固定的元模型，可有放回地抽取$K$个任务索引，携带每个任务的整组支持、查询与结果记录，重算任务等权均值。若元训练也被独立重复，任务与元训练种子是交叉的两层：分别抽任务索引和种子组索引，所有抽中的任务共用同一串种子索引，保留对应交叉记录。任务内支持集或查询对象另有抽样时，再按实际设计保留其分组，不能把任务、种子与照片摊平为独立样本。人工挑选的两个任务无法仅靠重抽样变成对广泛任务总体的代表；候选搜索的随机性也只有在重复整个选择过程时才被覆盖。

逐任务95\%区间不等于所有任务同时以95\%覆盖。若事先指定要同时判断$K$个任务均值，令总体允许的漏覆盖概率为$\alpha=0.05$，一种保守办法是Bonferroni校正：每项构造置信水平$1-\alpha/K$的双侧区间。在每项运行差满足上述$t$区间假设时，把式~\eqref{eq:ket-evaluation-paired-interval}~的分位数换成$t_{R-1,\,1-\alpha/(2K)}$即可。只要各项边际区间有相应覆盖率，全部区间同时覆盖的概率至少为$1-\alpha$，不要求任务间独立；代价是区间可能较宽。采用自助法时可相应改取$\alpha/(2K)$与$1-\alpha/(2K)$分位数，但覆盖率仍取决于重抽样近似，不能宣称有限样本精确保证。这个范围只包含预先列出的$K$项；若还比较多个方法、预算或事后挑选的任务，就要扩大比较族或另留确认数据。

任务顺序重复与初始化重复也不等价。固定顺序换种子只能说明该顺序下的稳定性；受控排列实验应报告哪些排列被采样、是否覆盖回归环境。真实时间流不能通过倒放时间来增加重复，可在不破坏时间的信息边界内重复初始化，并用不同时间窗口或不同实体群检验外推。可信性指标仍采用第一卷的攻击与删除验证口径，维护代价采用第~\ref{sec:ket-evaluation-resources}~节及式~\eqref{eq:ket-evaluation-ledger}~的资源定义，不能因为它们进入同一实验就改成普通准确率或参数量。

\section{评价结果与证据边界的分析}
\label{sec:ket-evaluation-analysis}

完整记录并不自动形成可靠结论。分析需要把收益分布、机制贡献和外推范围分开：前者说明谁受益，第二项判断哪些解释得到支持，第三项决定结论在什么条件下仍有根据。它们都应回到最初的主张，而不是只解释为何某个总分看起来较高。

\subsection{收益分布与能力取舍的分析}
\label{sec:ket-evaluation-distribution}

学习曲线和终点成绩应一起读。图~\ref{fig:ket-reuse-learning-curves}~中，迁移的固定标签优势从20标签处的17个百分点缩小到160标签处的1个百分点，早期优势没有消失，但幅度明显改变；第~\ref{sec:ket-reuse-prediction}~节另一组教学检查点则从早期领先转为后期落后。它们分别支持不同预算下的选择，不能用某一个早期点概括充分训练的效果。长序列还应把后续任务学习成本与历史任务保持并列，否则只看稳定旧分数会漏掉可塑性下降。

任务平均与样本平均也可能给出相反判断。设地区甲有900张照片，准确率从90\%变成92\%；地区乙有100张，从80\%变成70\%。这是教学构造，样本平均从$(810+80)/1000=89\%$变为$(828+70)/1000=89.8\%$，提高0.8个百分点；地区等权平均却从85\%降为81\%，降低4个百分点。前者按现有流量权重衡量总正确数，后者把两个地区视为同等重要。最差地区从80\%降为70\%，又揭示了另一项风险。应同时给出稀有类别、指定关键任务和预定权重的结果，不能事后挑一个有利平均。

即使效果相同，服务规模也会改变成本排序。图~\ref{fig:ket-evaluation-utility}~回顾第~\ref{chap:ket-reuse-selection}~章的教学账目：复用先支付准备费，曲线起点较高，但逐任务成本较低，增长较慢；从头训练没有这项准备费，却随任务数增长得更快。两线相交前后，成本较低的方案发生反转。图中只保留原例的同效果条件，实际判断还须把式~\eqref{eq:ket-evaluation-ledger}~中的维护、恢复与服务费用加回各方案。若另将准确率和计算合成效用，须给出权重与换算依据。

\begin{figure*}[htbp]
  \centering
  % Teaching costs from the reuse-selection chapter, equal target performance.
% C_reuse(m) = 120 + 10*m; C_scratch(m) = 30*m; all m = 1,...,10.
% Coordinate map: x = 24 + 12*m mm; y = 16 + 0.16*C mm.
\begin{tikzpicture}[
  x=1mm,y=1mm,font=\figurefont\footnotesize,text=BookInk,
  axis/.style={draw=BookInk,line width=0.7pt,-{Latex[length=1.8mm]}}
]
  \path[use as bounding box] (0,0) rectangle (169,82);
  \draw[BookTeal,line width=1pt] (29,76) -- (41,76);
  \fill[BookTeal] (35,76) circle (1mm);
  \node[anchor=west,inner sep=1mm] at (42,76) {复用};
  \draw[BookBlue,line width=1pt,dashed] (81,76) -- (93,76);
  \draw[BookBlue,fill=white,line width=0.8pt] (86,75) rectangle (88,77);
  \node[anchor=west,inner sep=1mm] at (94,76) {从头训练};

  \foreach \v in {0,100,200,300} {
    \draw[BookRule,line width=0.5pt] (24,{16+0.16*\v}) -- (144,{16+0.16*\v});
    \node[anchor=east,inner sep=1mm] at (23,{16+0.16*\v}) {$\v$};
  }
  \draw[axis] (24,16) -- (24,69);
  \draw[axis] (24,16) -- (151,16);
  \foreach \m in {0,1,2,4,6,8,10} {
    \draw[BookInk,line width=0.6pt] ({24+12*\m},16) -- ({24+12*\m},14.5);
    \node[anchor=north,inner sep=1mm] at ({24+12*\m},14.5) {$\m$};
  }
  \node[rotate=90] at (5,42) {总计算成本 $C$（GPU小时）};
  \node at (86,3) {实际服务任务数 $m$};

  \draw[BookTeal,line width=1pt]
    plot[domain=1:10,samples=10] ({24+12*\x},{16+0.16*(120+10*\x)});
  \draw[BookBlue,line width=1pt,dashed]
    plot[domain=1:10,samples=10] ({24+12*\x},{16+0.16*(30*\x)});
  \foreach \m in {1,...,10} {
    \fill[BookTeal] ({24+12*\m},{16+0.16*(120+10*\m)}) circle (1mm);
    \draw[BookBlue,fill=white,line width=0.8pt]
      ({23+12*\m},{15+0.16*(30*\m)})
      rectangle ({25+12*\m},{17+0.16*(30*\m)});
  }
  \draw[BookMuted,line width=0.7pt,dotted] (96,16) -- (96,44.8);
  \draw[BookMuted,line width=0.7pt] (96,46) -- (96,57);
  \node[anchor=south,inner sep=1mm] at (96,57) {$m=6$持平};
  \node[anchor=north west,inner sep=0pt,text=BookMuted] at (31,68)
    {成本越低越好};
\end{tikzpicture}

  \caption{复用规模改变成本排序的教学曲线，不是实测结果。两方案达到相同目标效果，推理成本相同，其他费用暂略；点来自$C_{\mathrm{reuse}}=120+10m$与$C_0=30m$（GPU小时）。解$120+10m=30m$得$m=6$时两者均为180；之后复用成本较低。横轴为实际服务任务数，连线只辅助观察整数点间趋势。}
  \label{fig:ket-evaluation-utility}
\end{figure*}

编辑结果沿第~\ref{sec:ket-evaluation-editing}~节的四组问题跟踪连续修改后的错误轨迹，模型整合则并列逐来源退化与组合留出的失败。发生取舍时，按第~\ref{chap:ket-sharing}~章约定的主辅关系或共同交付要求判断，明确哪些对象受益、哪些受损及其不确定性。

\subsection{机制贡献与替代解释的分析}
\label{sec:ket-evaluation-mechanisms}

观察到改善以后，应检查能否用更简单的条件差异解释。若迁移模型拥有更强预训练，收益可能来自预训练；若共享模型接触更多标注，收益可能来自额外监督；若更新方法多做若干轮训练，收益可能来自计算；若融合后保留更多参数和模型副本，收益可能来自容量。把完整方法、强基线和关键消融放在同一信息与预算表中，才能判断哪些解释已经被控制。

来源代理分数尤其需要与最终收益分别分析。高分来源是否在相同预算的实际适应中领先，目标试训是否改变排名，被淘汰来源是否只是早期收敛较慢，都是可检验的问题。排序相关性高，只说明候选间排序较一致；所有候选仍可能都不如无迁移对照。反过来，试训选择带来了更高终点效果，也可能同时耗尽额外标注和计算，应把选择后的收益与选择成本放回第~\ref{sec:ket-evaluation-target}~节的两种比较口径。

保持约束也可能用适应速度换取旧能力。可以固定记忆容量与新数据，把约束强度改变时的旧任务保持、新任务学习曲线和总计算并列；若同时换了骨干或增加回放样本，就不能单独解释为约束强度的作用。消融支持的是被实施的干预对完整系统的影响，只有干预确实隔离了目标因素，才有较强的机制解释。

机制假说、受控实验和实际环境证据承担不同职责。“任务梯度较一致时共享更有效”可以是待检验假说；固定其他条件改变任务搭配，可能提供受控支持；在新地区真实工作负载上复现，则补充外部有效性。仅在若干任务中看到梯度相似与收益相关，不能据此断言相似性导致收益。近期预印本也应按同样标准检查公开配置、重复、消融与跨数据复核，不以发表时间或框架名称决定可信程度；证据尚限于少数模型时，应保留这个范围。

\subsection{适用范围与证据缺口的分析}
\label{sec:ket-evaluation-limits}

结论至少应带上已检验的信息权限、任务分布、模型起点和预算。例如，“在提供类别描述、每任务五个支持样本的图像分类任务上，相对相同预训练的普通适应减少过线标签”不能缩写为“无需监督即可迁移”。如果未使用目标标签的来源选择只由代理分数支持，应说明还没有直接确认目标收益；第~\ref{chap:transfer-learning-theory}~章的不可识别性边界不会因为测了更多无标签分数而消失。

长序列保持和可塑性必须限定在实际序列长度、任务到达规则与容量范围内。三阶段矩阵能解释指标，不能证明无限任务下持续增长；同一类别空间上的真实时间流，也不等于类别无限增加的开放世界。事实编辑在已测关系链上保持一致，不代表所有隐含依赖都被更新；符号指令上的组合成功，也不能直接跨越到未测的跨模态或跨语言任务。整体研究问题和方法之间的关系可回看第~\ref{chap:ket-overview}~章，但概念上的共同性不替代每个范围的新证据。

“实验未发现失败”“在给定测试上达到标准”和“在明确假设下获得保证”是三种不同强度的陈述。有限测试没有发现泄漏，仍可能只是攻击不足；没有显著收益，可能是确实接近，也可能是任务或运行太少；相反，统计上可检测的极小差异未必具有实际用途。报告应同时给出变化幅度、估计不确定性、成功条件和未覆盖范围，而不是只给显著与否。

证据不足时，可以收窄主张，或围绕缺口安排新的评价。若疑问是目标领域覆盖不足，就补领域留出；若疑问是预算漏记，就补完整账目；若疑问是关联推理失败，就补成对关系链和必要前提。把多个访问条件、任务和分母都不相同的好分数拼接起来，并不能修补任何一个缺口。评价最终应逐项回答开篇的问题：新地区总体改善是否牺牲关键任务，历史能力相对哪个时刻得到保持，组合成功是否超出来源选择，以及这些收益付出了多少总成本。

\section{本章小结}
\label{sec:ket-evaluation-summary}

知识演进与迁移的收益需要由一组相互对应的记录支持：研究主张确定比较对象，指标规定参照和分母，数据与留出让所声称的变化接受检验，对照和访问协议排除条件不等，资源账本界定取得与维持收益的代价。少用标签、固定预算下更准确、历史能力保持、内容纠正和组合泛化可以同时出现，也可能互相取舍，不能由一个综合分数相互替代。

三任务矩阵表明，相对历史最佳的遗忘与相对刚学完的BWT可以给出不同信息；学前零样本收益又不同于开始学习后的成本改善。过时事实的纠正应按新的时间条件验证，不计为有害遗忘；机器遗忘则要对照未使用撤回记录的训练过程。共享和融合的评价保留逐任务、逐来源结果，真正的组合能力还要通过成分已见而组合未见的测试。

一份可用的结论不仅说明平均提高多少，还说明哪些任务受损、哪些运行未过线、选择与适应使用了什么信息，以及结论受哪些任务、时间和预算限制。这样的证据能支持下一轮来源选择、知识更新或迁移决策，同时保留在数据不足、收益不稳或成本不合算时不继续复用的依据。
